\documentclass[a4paper,11pt]{article}
\usepackage{lecturenotes}
\usepackage{graphicx}
\usepackage{enumitem}
\setlist{itemsep=2pt,parsep=0pt,topsep=4pt}
\lectureheader{COE718}{3}
\newcommand{\sourcefigure}[2]{%
  \par\smallskip\begin{center}
  \includegraphics[width=\linewidth,height=#2,keepaspectratio]{figures/#1}
  \end{center}\smallskip}
% Show the graphical field above each table; the table is transcribed below.
\newsavebox{\sourceimagebox}
\newlength{\sourcebottom}
\newcommand{\sourcefield}[2]{%
  \sbox{\sourceimagebox}{\includegraphics{figures/#1}}%
  \setlength{\sourcebottom}{#2\ht\sourceimagebox}%
  \begin{center}\includegraphics[width=\linewidth,trim=0 {\sourcebottom} 0 0,clip]{figures/#1}\end{center}}
\newcommand{\code}[1]{\texttt{#1}}
\begin{document}
\tableofcontents
\clearpage

\notesection{ARM Registers and Programming Model}
There are 15 general-purpose registers.
\begin{itemize}
  \item \textbf{Thumb mode:} R0--R7, general-purpose.
  \item \textbf{7 high registers:} R8--R12, only accessible with MOV, ADD, or CMP.
  \item \textbf{Stack pointer:} R13.
  \item \textbf{Link register:} R14.
  \item \textbf{Program counter:} R15.
\end{itemize}

\subnotesection{Data-processing Instructions}
\notebox{Of Note}{
RSB, \emph{Reverse Subtract}, \code{RSB R3, R1, R2}: $\rightarrow R_2-R_1\implies R_3$.

SNB, \code{SNB R3, R1, R2}: $\rightarrow R_1-R_2\implies R_3$.
}
\emph{Add info for `S'-suffix here.}

\subnotesection{Loading Constants}
\begin{description}
  \item[\code{MOV, Rd, constant}] Works for 0--255 and $-255$--$-1$.
  \item[\code{MVN, Rd, constant; rd < -constant}]
  Effectively doubles the number of constants. The assembler converts MOV with a negative constant to MVN.
  \item[\code{LDR}, \code{rd = constant}]
  Assembler pseudo-op, not an instruction. Converts to MOV or MVN if possible; otherwise converts to \code{LDR r\_d, [PC, \#imm]}.
\end{description}

\textbf{LDRH (Load Half Word)}

Loads a 16-bit memory half-word that is unmodified, and pads the remaining bits of the 32-bit register (B31--B16) with zeros.

\textbf{LDRSH (Load Signed Halfword)}

Interprets the 16 bits as signed (a two's-complement number), sign-extends it to 32 bits, and stores it into the register.

\textbf{STRH (Store Half Word)}

Stores the last 16 bits of a half-word into memory (i.e., \textbf{MSB 16 bits are NOT written into memory}).
\sourcefigure{image-105940.png}{0.40\textheight}

\subnotesection{Condition Code Field}
\sourcefield{image-110008.png}{0.62}
The condition code is the 4-bit field $c_3\,c_2\,c_1\,c_0$.
\begingroup\small
\begin{notetable}{@{}lXlX@{}}{\textbf{cond} & \textbf{Meaning} & \textbf{cond} & \textbf{Meaning}}
0000 (0x0) & EQ: $Z=1$ & 1000 (0x8) & HI: $C=1$ and $Z=0$ \\
0001 (0x1) & NE: $Z=0$ & 1001 (0x9) & LS: $C=0$ or $Z=1$ \\
0010 (0x2) & CS/HS: $C=1$ & 1010 (0xA) & GE: $N=V$ \\
0011 (0x3) & CC/LO: $C=0$ & 1011 (0xB) & LT: $N\ne V$ \\
0100 (0x4) & MI: $N=1$ & 1100 (0xC) & GT: $Z=0$ and $N=V$ \\
0101 (0x5) & PL: $N=0$ & 1101 (0xD) & LE: $Z=1$ or $N\ne V$ \\
0110 (0x6) & VS: $V=1$ & 1110 (0xE) & AL: always \\
0111 (0x7) & VC: $V=0$ & 1111 (0xF) & NV: never / unused \\
\end{notetable}
\endgroup

\Needspace{18\baselineskip}
\textbf{Branch Instructions}
\begingroup\small
\begin{notetable}{@{}>{\raggedright\arraybackslash}p{0.21\linewidth}>{\raggedright\arraybackslash}p{0.28\linewidth}c>{\raggedright\arraybackslash}X@{}}{\textbf{Instruction} & \textbf{Operation} & \textbf{\{S\}} & \textbf{Notes}}
\code{B \{c\} Label} & $\mathrm{PC}\leftarrow\mathrm{PC}+\mathrm{imm}$ & n/a & `c' is an optional condition code. \\
\code{BL Label} & $\mathrm{PC}\leftarrow\mathrm{PC}+\mathrm{imm}$\newline $\mathrm{LR}\leftarrow\text{rtn adr}$ & n/a & Sub-routine call. \\
\code{BX Reg} & $\mathrm{PC}\leftarrow\mathrm{reg}$ & n/a & \code{BX LR}: return from function/procedure. \\
\code{CBZ R\_n, label} & If $R_n=0$,\newline $\mathrm{PC}\leftarrow\mathrm{PC}+\mathrm{imm}$ & n/a & Cannot append a condition code to CBZ. \\
\code{CBNZ R\_n, label} & If $R_n\ne0$,\newline $\mathrm{PC}\leftarrow\mathrm{PC}+\mathrm{imm}$ & n/a & Cannot append a condition code to CBZ. \\
$\mathrm{IT}_{C1C2C3}$ cond & Each $c_i$ is T, E, or EMPTY & n/a & Controls 1--4 instructions in an \textbf{IT} block. \\
\end{notetable}
\endgroup

\notesection{xPSR}
See Lecture 2, \emph{Embedded CPUs ARM7 and M3}, ``PSR Program Status Register'' (link text: \emph{ARM Interrupt Processing}).

The xPSR is a group of special registers that indicate the status of the program. It consists of the APSR, EPSR, and IPSR. The \code{GE[3:0]} is only available for the Cortex-M3 and some other ARMv6-M chips, and indicates the Greater-Than or Equal flags for each byte lane.

Thumb state is always 1.
\sourcefield{image-110803.png}{0.68}
\begingroup\small
\begin{notetable}{@{}>{\raggedright\arraybackslash}p{0.16\linewidth}>{\raggedright\arraybackslash}p{0.25\linewidth}>{\raggedright\arraybackslash}X@{}}{\textbf{Bit(s)} & \textbf{Name} & \textbf{Description}}
N & Negative flag & Result is negative (sign bit = 1). \\
Z & Zero flag & Result is zero. \\
C & Carry (NOT borrow) flag & Carry out (or NOT borrow for subtraction). \\
V & Overflow flag & Signed overflow occurred. \\
Q & Sticky saturation flag & Set if a saturation (Q) occurred (not available in ARMv6-M). \\
GE[3:0] & Greater-Than or Equal flags & Per-byte GE flags for SIMD/parallel operations (ARMv7E-M only; not available in ARMv6-M or Cortex-M3). \\
ICI/IT & Interrupt-Continuable Instruction & IC bits, IF-THEN instruction status for conditional execution (not available in ARMv6-M). \\
T & Thumb state & Always 1 in Thumb state; clearing this bit causes a fault exception. \\
Exception Number & Exception Number & Indicates which exception the processor is handling. \\
\end{notetable}
\endgroup

\Needspace{23\baselineskip}
\subnotesection{SSAT or USAT Instruction}
See Lecture 2, \emph{Embedded CPUs ARM7 and M3}, ``Saturation'' (link text: \emph{Saturation Inst.}).

Typically organized as \code{op\{cond\} R\_d, \#n, Rm \{, shift \#s\}}.
\begin{itemize}
  \item \textbf{op (SSAT/USAT):} Saturates a signed value to a signed range, or a signed value to an unsigned range.
  \item \textbf{Cond:} Condition code.
  \item \textbf{Rd:} Specifies the destination register.
  \item \textbf{n:} Specifies the bit position to saturate to (0--32 for SSAT and 0--31 for USAT).
  \item \textbf{Rm:} Register containing the value to saturate.
\end{itemize}
\[
\begin{aligned}
\mathrm{SSAT}:&\quad 2^{n-1}\le x\le 2^{n-1}-1,\\
\mathrm{USAT}:&\quad 0\le x\le 2^n-1.
\end{aligned}
\]
If the returned result differs from the value to be saturated, it is called \emph{saturation}. If this occurs, the Q flag is set to 1.

Example: \code{SSAT R7, \#16, R7, LSL, \#4}.

\notesection{Bit-banding}
You

The M3 has memory-mapped I/O, with 4 GB of memory space organized in bytes, which is quite large. Bit-banding takes advantage of this by using two different memory regions to refer to the same physical data.
\begin{itemize}
  \item \textbf{Primary Bit Band Region:} Each address corresponds to a single data byte.
  \item \textbf{Bit Band Alias:} Each address corresponds to 1 bit of the same data.
\end{itemize}
This allows access to a bit of data by a single instruction.

Bit-banding does not allow interruption of read-modify-write.

\textbf{Applications}
\begin{itemize}
  \item Access individual status and control bits of an I/O device, or implement a set of 1-bit Boolean flags (\emph{which can be used to implement Mutex objects}).
\end{itemize}
\keyequation{Bit-band Word Address}{
\begin{aligned}
\text{Bit-Band Word Address} ={}& \text{Bit Band Alias Base Address}\\
 &+ (\text{Byte Offset}\times32) + (\text{Bit Number}\times4)
\end{aligned}}
Where \emph{Byte Offset} is:
\[
\text{Byte Offset}=\text{Bit's Bit Band Address}-\text{Bit Band Base Address}.
\]
The address range \code{0x20000000} is SRAM, and the peripheral space is \code{0x40000000}.

\Needspace{7\baselineskip}
Bit-banding is for two predefined memory regions:
\begin{multicols}{2}
\begin{itemize}
  \item First 1 MB of SRAM.
  \item First 1 MB of the Peripheral Region.
\end{itemize}
\end{multicols}
One bit is addressed by its own 32-bit in a separate part of memory.

\Needspace{8\baselineskip}
\textbf{General Steps for Bit-banding}
\begin{enumerate}
  \item Calculate word address.
  \item Define a pointer to the address: \code{\#define BIT\_ADDR = (*(...*))}.
  \item Assign a value to the port bit.
\end{enumerate}

\notesection{Conditional Execution}
Conditional execution is typically a standard ALU instruction, i.e., ADD, with a conditional statement added as a suffix: \code{ADDEQ}. This instruction will only execute when the zero flag in the CPSR is set.

\notesection{IT Block (If-Then)}
The IT block makes up the four following instructions. Its general form is:
\begin{notecode}[language={}]
IT {x {y {z} } } {cond}
\end{notecode}
\begin{itemize}
  \item \emph{x} specifies the condition switch for the second instruction in an IT block.
  \item \emph{y} specifies the condition switch for the third instruction in the IT block.
  \item \emph{z} specifies the condition switch for the fourth instruction in the IT block.
  \item \emph{cond} specifies the condition for the first instruction in the IT block.
\end{itemize}
The condition switch for the 2nd, 3rd, and 4th instruction in the IT block is either:
\begin{itemize}
  \item \emph{T}, Then: Applies the condition cond to the instruction.
  \item \emph{E}, Else: Applies the inverse condition of cond to the instruction.
\end{itemize}
The instructions (including branches) in the IT block, except the \emph{BKPT} instruction, must specify the condition in the \code{\{cond\}} part of their syntax.

The assembler auto-generates the IT blocks according to the conditions with respect to instructions.
\begin{itemize}
  \item With the exception of CMP, CMN, and TST, the 16-bit instructions that normally affect the condition code flags do not affect them in the IT block.
  \item A BKPT instruction is always executed; it does not need a condition in the \code{\{cond\}} part of the syntax.
  \item Conditional branches have a longer branch range than outside the IT block.
\end{itemize}
The following instructions are not allowed inside an IT block: \code{IT, CBZ \& CBNZ, TBB \& TBH, CPS \& CPSID \& CPSIE, SETEND}.

\textbf{Some other restrictions:}
\begin{itemize}
  \item A branch or any instruction that modifies PC is only allowed as the last instruction in the block.
  \item Cannot branch to any instruction in an IT block, unless returning from an exception handler.
\end{itemize}
\Needspace{4\baselineskip}
This instruction is only available as a 16-bit instruction in ARMv6T2 and above, and is pseudocode in regular ARM code.

ARM allows non-control-flow-based instructions to be appended with conditional codes, allowing more efficient coding and processor performance:
\begin{notecode}[language=ARM]
CMP r2, #5 //if(a<=5);
MOVELE R2, #10 //a = 10;
MOVGT r2, #1 //else a=1;
\end{notecode}

\subnotesection{ITTE Block}
\sourcefigure{image-004555.png}{0.48\textheight}

\Needspace{25\baselineskip}
\notesection{Loop Variables}
\begin{notecode}[language=C]
while (a!= b){
    if(a>b) a = a-b;
    else b = b-a;
}
\end{notecode}
\begin{notecode}[language=ARM]
LDR R0, a
LDR R1, b
top: CMP R0, R1
     BEQ done
     ITE GT
     SUBGT R0, R0, R1
     SUBLE R1, R1, R0
     B top
done:
     ; R0 = R1 = GCD(a,b)
\end{notecode}

\notesection{Function Call and Return}
\textbf{Function Call:} \code{BL Function}
\begin{itemize}
  \item Loads the program counter (PC) with the entry-point address of the function.
  \item Saves the return address in the link register.
\end{itemize}
\textbf{Function Return:} \code{BX lr}
\begin{itemize}
  \item Copies the link register back into the program counter.
\end{itemize}

\Needspace{0.65\textheight}
\notesection{Temporary Variables}
\sourcefigure{image-015921.png}{0.53\textheight}
\emph{r0--r3} are not used for any parameters and must be preserved and restored around any call.\par
\emph{r4--r8} must always be preserved and restored.

\textbf{Essentially:}
\begin{enumerate}
  \item \code{r0-r3} are temporary (scratch) registers: functions can freely use/overwrite them for parameters, return values, and temporary variables.
  \item \code{r4-r8} must be preserved. If a function uses them, it must \textbf{push} their original values onto the stack and \textbf{pop} them before returning.
  \item LR stores the return address. If a function calls another function with BL, it should save LR first because the new BL will overwrite it.
\end{enumerate}

\Needspace{14\baselineskip}
\notesection{Cortex A, X, and C1 Series CPU}
The \textbf{A72} has $3.5\times$ the performance of the A15, and is paired with the A53 to get \emph{big little} architecture.

\textbf{Cortex A75}
\begin{itemize}
  \item Can execute 3 instructions per clock cycle and has 7 execution units.
  \item Has two load/stores, two NEON \& FPUs, a BPU, and two integer cores.
\end{itemize}

\subnotesection{Unified C1 Family}
\begingroup\small
\begin{notetable}{@{}>{\raggedright\arraybackslash}p{0.16\linewidth}>{\raggedright\arraybackslash}p{0.17\linewidth}>{\raggedright\arraybackslash}X>{\raggedright\arraybackslash}X@{}}{\textbf{Core Name} & \textbf{Predecessor} & \textbf{Key Benefit} & \textbf{Ideal Uses}}
Cortex-X925\newline\textcolor{red}{(2024)} & Cortex-X4\newline\textcolor{red}{(2023)} & Ultimate performance (+36\% vs X4) & Flagship smartphones, laptops \\
C1-Ultra\newline\textcolor{red}{(2025)} & Cortex-X925 & Flagship peak performance (+25\% vs X925) & Generative AI, content creation \\
C1-Premium & New Category & Ultra performance, 35\% smaller area & Sub-flagship mobile, multitasking \\
C1-Pro & Cortex-A725 & +16\% sustained performance & Video playback, gaming \\
C1-Nano & Cortex-A520 & 26\% more power-efficient & Wearables (smart watch), background tasks \\
\end{notetable}
\endgroup
The smartphone cores are classified across four categories: Lumex C1-Ultra, Lumex C1-Premium, Lumex-C1 Pro, and the Lumex C1-Nano.

\subnotesection{ARMv7 Suffixes}
Instructions do not update the PSR unless a suffix \textbf{S} is appended to the instruction (exceptions: \emph{CMP} and \emph{Test}, e.g., TST, TEQ, which write to the PSR directly).

Consists of 16-bit Thumb instructions.

There also exist suffixes to indicate the size of the word to be loaded or stored, in the form \emph{LDR(size).W}.

\subnotesection{Conditional Execution}
It exists, and executes individual instructions based on the condition flags set by the previous instruction(s). It can be invoked by using branches or adding a condition-code suffix to the instruction.

\textbf{Table 4.1: Suffixes in Instructions}
\begin{notetable}{@{}>{\raggedright\arraybackslash}p{0.22\linewidth}>{\raggedright\arraybackslash}X@{}}{\textbf{Suffix} & \textbf{Description}}
S & Update Application Program Status register (APSR) (flags); for example: \code{ADDS R0, R1 ; this will update APSR}. \\
EQ, NE, LT, GT, and so on & Conditional execution; EQ = Equal, NE = Not Equal, LT = Less Than, GT = Greater Than, and so forth. For example: \code{BEQ <label> ; branch if equal}. \\
\end{notetable}
\Needspace{20\baselineskip}
\textbf{BNE, BLT, BGT, etc.}
\begingroup\small
\begin{notetable}{@{}l>{\raggedright\arraybackslash}X@{}}{\textbf{Instruction} & \textbf{Description}}
B & Branch. \\
B<cond> & Conditional branch. \\
BL & Branch with link; call a subroutine and store the return address in LR (this is actually a 32-bit instruction, but it is also available in Thumb in traditional ARM processors). \\
BLX & Branch with link and change state (BLX <reg> only)\textsuperscript{1}. \\
BX <reg> & Branch with exchange state. \\
CBZ & Compare and branch if zero (architecture v7). \\
CBNZ & Compare and branch if nonzero (architecture v7). \\
IT & IF-THEN (architecture v7). \\
\end{notetable}
\endgroup
IT blocks handle small conditional code and are used to avoid branch penalties. They have a maximum of four conditionally executed instructions.

\subnotesection{Bit-banding}
Address \code{0x20000000}: SRAM.\par
Address \code{0x40000000}: Peripheral.
\[
\begin{aligned}
\text{Bit band word address} ={}& \text{Bit band alias Base address}\\
 &+ (\text{Byte Offset}\times32) + (\text{Bit number}\times4).
\end{aligned}
\]
Where:
\[
\text{byte offset}=\text{Bit's Bit band base address}-\text{Bit band base address}.
\]
Bit-banding enables fewer instructions for bit manipulation/access.
\end{document}
